%%%%%%
% Metadata
%%%%%%
\chapter{Testing}

\section{Study Selection for Evaluation}


This section describes the source of the studies used to evaluate the proposed alpha recommendation system and explains the criteria used for their selection.

The studies used in this thesis are drawn from the FORRT Replication Database (FReD), a publicly available database that compiles published replication studies across multiple research fields. These studies are used to examine whether the recommended alpha values generated by the system would have led to more reliable statistical decisions. Each entry in the database links an original study with its corresponding replication outcome and includes key study characteristics such as sample size and effect size.

The database contains studies from several disciplines, including psychology, behavioral science, social science, economics, and education. These fields were selected because null hypothesis significance testing is widely used and documented replication results are available.

For this evaluation, 20 studies were selected based on three criteria. First, the original study reported a statistically significant result using the conventional threshold of $\alpha = 0.05$. Second, an independent replication study was conducted. Third, the replication did not support the original finding. These cases are important because they show situations where a result was initially considered significant but later failed to replicate. By focusing on these studies, the evaluation examines whether a different choice of alpha, as recommended by the system, could have led to a more cautious and potentially more reliable decision.

The selected studies differ in sample size, effect size, and research field. This variation allows the system to be tested across different study settings rather than within a single discipline.

\begin{table}[ht]
	\centering
	\caption{Selected Studies and Recommended Alpha Values}
	\label{tab:selected_studies}
	
	\resizebox{\textwidth}{!}{
		\begin{tabular}{cccccc}
			\hline
			Study & Field & Sample Size & Effect Size & Replication Outcome & Recommended Alpha \\
			\hline
			1  & Psychology & 33  & 0.43 & Failed & 0.0325 \\
			2  & Psychology & 33  & 0.62 & Failed & 0.063  \\
			3  & Psychology & 34  & 0.54 & Failed & 0.050  \\
			4  & Psychology & 35  & 0.61 & Failed & 0.060  \\
			5  & Psychology & 36  & 0.54 & Failed & 0.042  \\
			6  & Behavioral Science & 170 & 0.19 & Failed & 0.040 \\
			7  & Behavioral Science & 170 & 0.19 & Failed & 0.040 \\
			8  & Behavioral Science & 172 & 0.15 & Failed & 0.050 \\
			9  & Behavioral Science & 172 & 0.15 & Failed & 0.050 \\
			10 & Behavioral Science & 139 & 0.17 & Failed & 0.0537 \\
			11 & Behavioral Science & 424 & 0.10 & Failed & 0.020 \\
			12 & Behavioral Science & 424 & 0.11 & Failed & 0.020 \\
			13 & Behavioral Science & 398 & 0.12 & Failed & 0.030 \\
			14 & Behavioral Science & 398 & 0.18 & Failed & 0.020 \\
			15 & Psychology & 48  & 0.40 & Failed & 0.030 \\
			16 & Psychology & 168 & 0.25 & Failed & 0.040 \\
			17 & Psychology & 41  & 0.31 & Failed & 0.020 \\
			18 & Psychology & 85  & 0.25 & Failed & 0.040 \\
			19 & Education & 6   & 0.52 & Failed & 0.080 \\
			20 & Education & 744 & 0.27 & Failed & 0.001 \\
			\hline
		\end{tabular}
	}
	
\end{table}

\section{Testing the Developed Alpha Recommendation Approach}

This section evaluates the developed alpha recommendation system using the 20 studies identified in Section~4.

All selected studies originally rejected the null hypothesis using the conventional significance level of $\alpha = 0.05$. However, these findings failed to replicate. These cases therefore represent situations in which the original statistical decision was likely a false positive.

For each study, the required inputs—sample size and effect size—were entered into the alpha recommendation system. Based on these inputs, the system generated a study-specific recommended alpha value. These recommended values were then compared with the fixed threshold of $\alpha = 0.05$.

\subsection{Decision Comparison}

The recommended alpha values ranged from 0.001 to 0.08, demonstrating that the system adapts the significance level according to study characteristics rather than applying a single fixed threshold.

Among the 20 studies, 11 received a recommended alpha value lower than 0.05, indicating a stricter decision rule. Three studies received a recommendation equal to 0.05, and six received a recommendation higher than 0.05.

In the 11 cases where the recommended alpha was stricter than 0.05, the null hypothesis would not have been rejected under the proposed system. Because these studies later failed to replicate, these cases represent instances in which the system would have avoided potentially misleading rejections.

In the three cases where the recommended alpha matched 0.05, the system’s recommendation was consistent with the conventional threshold.

In the six cases where a more lenient alpha was recommended, the system reflected a higher tolerance for false positives based on study characteristics. A more lenient alpha does not imply that the original result was correct; rather, it reflects the contextual balance between Type I and Type II errors within that specific study.

\subsection{Overall Result}

Compared with the routine use of $\alpha = 0.05$, the proposed alpha recommendation system reduced the number of potentially misleading rejections in many failed replication cases. These findings suggest that treating alpha as a context-dependent quantity can improve statistical decision-making.

\section{Conclusion}

This thesis proposed a data-driven approach for selecting the significance level in hypothesis testing. Instead of relying on a fixed threshold such as $\alpha = 0.05$, the proposed system adapts the significance level based on study characteristics.

Using 20 failed replication studies from multiple research fields, the evaluation demonstrated that the recommended alpha values often differed from 0.05 and, in many cases, would have avoided misleading rejections of the null hypothesis. These results support the argument that alpha should not be treated as a universal constant.

Overall, the findings indicate that a flexible and context-sensitive approach to selecting alpha can improve statistical decision-making across research domains.